\documentclass[twoside,11pt]{report}
\usepackage{titlesec}
\usepackage{multirow}
\usepackage{subfig}
%\usepackage{wasysym}
\usepackage{tikz-cd}
\usetikzlibrary{arrows}
\usetikzlibrary{cd}
\usepackage{bbm}
\usepackage{rotating}
%\usetikzlibrary{automata,positioning}
\newcommand{\mychapter}[2]{
    \setcounter{chapter}{#1}
    \setcounter{section}{0}
    \chapter*{#2}
    \addcontentsline{toc}{chapter}{#2}
}


\usepackage{estilo-apuntes}
\usepackage{braids}
\usepackage{pgfplots}
\pgfplotsset{compat=1.13}
\usetikzlibrary{decorations.shapes}
\pgfkeyssetvalue{/tikz/braid height}{1cm} %no parece hacer nada
\pgfkeyssetvalue{/tikz/braid width}{1cm}
\pgfkeyssetvalue{/tikz/braid start}{(0,0)}
\pgfkeyssetvalue{/tikz/braid colour}{black}

\bibliographystyle{amsalpha} %otros: ieeetr, IEEEtran

\begin{document}
\begin{titlepage}
	\centering
		\includegraphics[width=8cm]{Imagenes/sello.png}
		
	{\Large\bfseries UNIVERSIDAD DE SEVILLA\par}
	{\Large\bfseries FACULTAD DE MATEMÁTICAS\par}
	\vspace{0.5cm}
	{\large\bfseries Trabajo de Fin de Máster \par}
	\vspace{1cm}
	{\Huge\bfseries La conjetura de Deligne\par}
	\vspace{1cm}
	{\large Por:\par}
	{\large Javier Aguilar Martín\par}
	\vspace{0.5cm}
	{\large Dirigido por:\par}
	{\large Fernando Muro Jiménez\par}
	\vspace{1cm}
	{\large Máster Universitario en Matemáticas\par}
	\vspace{0.5cm}
	{\large  Junio de 2019\par}
	\vspace{1cm}





	
\end{titlepage}



\tableofcontents

%\url{https://tex.stackexchange.com/questions/117745/bibtex-display-url-in-reference} PARA PONER LAS URL EN LA BIBLIOGRAFÍA QUE NO SALEN

%\url{https://tex.stackexchange.com/questions/119805/bibliography-in-texmaker}\url{https://tex.stackexchange.com/questions/8458/making-the-bibliography-appear-in-the-table-of-contents}

%EJERCICIO 
%
%Sean $(X,A)$ y $(Y,B)$ pares que satisfacen la HEP (se pueden considerar subcomplejos si es necesario) tales que $A$ y $B$ son contráctiles. Sea $f:(X,A)\to(Y,B)$ una aplicación de pares. Sean $p_X:\widetilde{X}\to X$ y $p_Y:\widetilde{Y}\to Y$ aplicaciones recubridoras. Probar que $p^{-1}(A)$ es unión de componentes homeomorfas a $A$ a través de $p$ y lo mismo para $B$. Si $\widetilde{A}$ y $\widetilde{B}$ son unas de estas componentes, entonces probar que existe una única $\tilde{f}$ que hace conmutar el diagrama
%\[
%\begin{tikzcd}
%(\widetilde{X},\widetilde{A})\arrow[r, "\tilde{f}"]\arrow[d, "p_X"] & (\widetilde{Y}, \widetilde{B})\\
%(X,A)\arrow[r, "f"] & (Y,B)
%\end{tikzcd}
%\]

\subfile{Abstract}
\subfile{Resumen} %NO ESTOY SEGURO DE SI HACE FALTA
\subfile{Intro} %SI PONGO BASTANTE EN EL ABSTRACT NO HARÍA FALTA
\subfile{Gerstenhaber}
\subfile{Nerve}
\subfile{Operads}
\subfile{Formality}
\subfile{Proof}
\subfile{Bibliografia}

%¿AÑADO ALGUNA COSA MÁS DE ESTILO?




%Si me preguntan sobre la característica 0: las cosas de teoría de representaciones y que en char 2 tienes que decir que el corchete cumple [x,x]=0 porque con antisimetría es trivial






\end{document}
